\documentclass[11pt]{article}
\usepackage[a4paper,margin=28mm]{geometry}
\usepackage{amsmath,amssymb,amsthm}
\usepackage[hidelinks]{hyperref}
\newtheorem{theorem}{Theorem}
\newtheorem{lemma}[theorem]{Lemma}
\newtheorem{corollary}[theorem]{Corollary}
\theoremstyle{remark}
\newtheorem{remark}[theorem]{Remark}
\newcommand{\gen}[1]{\langle #1\rangle}
\newcommand{\Stab}{\operatorname{Stab}}
\newcommand{\GL}{\operatorname{GL}}
\newcommand{\PGL}{\operatorname{PGL}}
\title{A stabilizer certificate for the intersection property in rank four}
\author{}
\date{}
\begin{document}
\maketitle

\noindent
The purpose of this note is to state precisely a sufficient certificate for
the intersection of two end subgroups and to give a complete proof that the
certificate implies the rank-four intersection property. The ingredients are
classical: Schreier's stabilizer generators and the end-subgroup criterion
for string C-groups. No new general theorem, classification result, successful
finite computation or formal verification is claimed here. The arguments are
included in full so that the conditions required of a finite certificate are
unambiguous.

\section{Groups and actions}

Let $G=\gen{a,b,c,e}$, where $a,b,c,e$ are elements of order exactly two.
Assume the string commutation relations
\begin{equation}\label{eq:string}
 [a,c]=[a,e]=[b,e]=1.
\end{equation}
Put
\[
 L=\gen{a,b,c},\qquad R=\gen{b,c,e},\qquad D=\gen{b,c}.
\]
All subgroup equalities in this note take place in the same group $G$.
For a distinguished list $(r_i)_{i\in M}$, write
$G_I=\gen{r_i:i\in I}$ for $I\subseteq M$, with $G_\varnothing=1$.
Its intersection property is
\begin{equation}\label{eq:IP}
 G_I\cap G_J=G_{I\cap J}\qquad(I,J\subseteq M).
\end{equation}
A string C-group is a group generated by its distinguished involutions,
satisfying the string commutation relations and \eqref{eq:IP}.

An action on a set $\Omega$ is written on the right:
\[
 (x^g)^h=x^{gh},\qquad x^1=x.
\]
No faithfulness hypothesis is imposed on this action. For $x\in\Omega$,
$\Stab_L(x)$ means its stabilizer inside $L$, not inside a larger ambient
permutation group. An element fixes a projective point when it preserves the
corresponding line; it need not fix a chosen nonzero vector on that line.

\section{The stabilizer certificate}

\begin{lemma}[Schreier generators in the chosen convention]\label{lem:schreier}
Let $L=\gen{S}$ act on the right on $\Omega$, and assume $S=S^{-1}$.
Fix $\ell\in\Omega$. Suppose $\mathcal O\subseteq\Omega$ is a finite set
containing $\ell$, closed under every $s\in S$. For each $x\in\mathcal O$,
suppose an element $T_x\in L$ is specified such that
\[
 \ell^{T_x}=x,\qquad T_\ell=1.
\]
Then $\mathcal O=\ell^L$, and
\begin{equation}\label{eq:schreier}
 \Stab_L(\ell)=
 \gen{T_x s T_{x^s}^{-1}:x\in\mathcal O,\ s\in S}.
\end{equation}
In particular, the representatives need not form a prefix-closed transversal.
\end{lemma}

\begin{proof}
Every word in $S$ maps $\ell$ into $\mathcal O$, by closure, so
$\ell^L\subseteq\mathcal O$. The specified representative $T_x$ proves
$x\in\ell^L$ for every $x\in\mathcal O$, giving the reverse inclusion.

For $y=x^s$, the action convention gives
\[
 \ell^{T_xsT_y^{-1}}=(x^s)^{T_y^{-1}}
 =y^{T_y^{-1}}=\ell.
\]
Consequently every generator on the right side of \eqref{eq:schreier}
belongs to $\Stab_L(\ell)$.

Conversely, take $h\in\Stab_L(\ell)$ and write
$h=s_1\cdots s_m$ with $s_i\in S$.
Set $x_i=\ell^{s_1\cdots s_i}$ and $t_i=T_{x_i}$, including
$x_0=\ell$ and $t_0=1$. Since $h$ fixes $\ell$, also $x_m=\ell$ and
$t_m=1$. Thus
\[
 k_i=t_{i-1}s_i t_i^{-1}
\]
is one of the displayed Schreier generators, and direct cancellation gives
\[
 k_1\cdots k_m
 =t_0s_1t_1^{-1}t_1s_2t_2^{-1}\cdots t_{m-1}s_mt_m^{-1}
 =s_1\cdots s_m=h.
\]
The identity element corresponds to the empty word. This proves the claimed
generation and does not use faithfulness or finiteness of $L$.
\end{proof}

\begin{theorem}[Certificate for the end intersection]\label{thm:ends}
Let $L=\gen{a,b,c}$, $R=\gen{b,c,e}$ and $D=\gen{b,c}$ be as above,
acting on the right on $\Omega$. Suppose that:
\begin{enumerate}
\item $R$ fixes a point $\ell\in\Omega$;
\item a finite set $\mathcal O$ and representatives $T_x$ satisfy all the
      hypotheses of Lemma~\ref{lem:schreier} for $S=\{a,b,c\}$;
\item for every $x\in\mathcal O$ and $g\in\{a,b,c\}$, the exact group
      element $T_xgT_{x^g}^{-1}$ belongs to $D$.
\end{enumerate}
Then
\[
 \Stab_L(\ell)=D\qquad\hbox{and}\qquad L\cap R=D.
\]
\end{theorem}

\begin{proof}
The generators $a,b,c$ are involutions, so they form an inverse-closed
generating set. Lemma~\ref{lem:schreier} and the third assumption give
$\Stab_L(\ell)\leq D$. Since $D\leq L\cap R$ and $R$ fixes $\ell$,
\[
 D\leq L\cap R\leq\Stab_L(\ell)\leq D.
\]
All inclusions are therefore equalities.
\end{proof}

It suffices to verify that $b,c,e$ fix $\ell$; their generated group $R$
then fixes it. A complete finite table of $D$ permits the last assumption
to be checked by membership in that table. The theorem does not require
enumeration of $L$, $R$, or $G$, or knowledge of their orders. It is a
sufficient method: failure to find such a point, or a stabilizer larger than
$D$, does not itself disprove $L\cap R=D$.

\section{A complete proof of the rank-four criterion}

\begin{theorem}[The classical end criterion in rank four]\label{thm:rank4}
Assume \eqref{eq:string}. Suppose $(L;(a,b,c))$ and $(R;(b,c,e))$
satisfy the intersection property, and suppose $L\cap R=D$.
Then $(G;(a,b,c,e))$ is a string C-group of rank four.
\end{theorem}

\begin{proof}
For clarity, put
\[
 A=\gen a,\quad B=\gen b,\quad C=\gen c,\quad E=\gen e,
 \quad P=\gen{a,b},\quad Q=\gen{c,e}.
\]
The intersection properties in $L$ and $R$ first give
\begin{equation}\label{eq:basic}
 L\cap E=1,\qquad R\cap A=1,\qquad P\cap R=B,
 \qquad L\cap Q=C,
 \qquad P\cap Q=1.
\end{equation}
Here are the details. Since $E\leq R$, we have
$L\cap E=(L\cap R)\cap E=D\cap E=1$, using the intersection property
inside $R$; similarly $R\cap A=D\cap A=1$ inside $L$.
Since $P\leq L$, the equality
$P\cap R=P\cap D=B$ follows inside $L$.
Likewise, $L\cap Q=D\cap Q=C$ inside $R$.
Finally $P\cap Q\leq P\cap R=B$, and $B\cap Q=1$ inside $R$.

The two remaining maximal distinguished subgroups are
\[
 X=\gen{a,b,e},\qquad Y=\gen{a,c,e}.
\]
By \eqref{eq:string}, $e$ commutes with every element of $P$, and
$a$ commutes with every element of $Q$. Consequently
\[
 X=PE,\qquad Y=AQ.
\]
The first two equalities in \eqref{eq:basic} imply
$P\cap E=1$ and $A\cap Q=1$. Thus these products are internal direct
products: every element of $X$ has a unique expression $pf$ with
$p\in P,f\in E$, and every element of $Y$ has a unique expression $tq$
with $t\in A,q\in Q$.

Both $X$ and $Y$ have the intersection property on their distinguished
generators. To see this explicitly for $X$, for
$I\subseteq\{a,b,e\}$ let
$P_I=\gen{I\cap\{a,b\}}$ and $E_I=\gen{I\cap\{e\}}$.
Then $\gen I=P_IE_I$, and uniqueness of the direct-product coordinates
implies
\[
 (P_IE_I)\cap(P_JE_J)=(P_I\cap P_J)(E_I\cap E_J).
\]
The group $P$ inherits the intersection property from $L$, and the group
$E$, with its single distinguished involution, has that property directly.
Hence the displayed expression is $P_{I\cap J}E_{I\cap J}$.
This proves the assertion for $X$. The same argument, using $A$ and the
subgroup $Q$ of $R$, proves it for $Y$.

We now establish all six intersections of the maximal distinguished
subgroups:
\begin{equation}\label{eq:six}
\begin{array}{lll}
 L\cap R=D,&L\cap X=P,&L\cap Y=AC,\\[2pt]
 R\cap X=BE,&R\cap Y=Q,&X\cap Y=AE.
\end{array}
\end{equation}
The first is a hypothesis. If $x=pf\in L\cap X$ with $p\in P$ and
$f\in E$, then $f=p^{-1}x\in L\cap E=1$, giving $L\cap X=P$.
If $x=tq\in R\cap Y$ with $t\in A$ and $q\in Q$, then
$t=xq^{-1}\in R\cap A=1$, giving $R\cap Y=Q$.
For $x=tq\in L\cap Y$, we instead have $q=t^{-1}x\in L\cap Q=C$,
so $x\in AC$; the reverse inclusion is immediate.
For $x=pf\in R\cap X$, we have $p=xf^{-1}\in P\cap R=B$,
so $x\in BE$, again with the reverse inclusion immediate.

It remains to justify $X\cap Y=AE$. Take $x\in X\cap Y$ and write
$x=pf=tq$ as above. Then
\[
 t^{-1}p=qf^{-1}\in P\cap Q=1,
\]
since $A\leq P$ and $E\leq Q$. Thus $p=t\in A$ and $q=f\in E$,
which gives $x\in AE$. Conversely $A,E\leq X\cap Y$, so equality
holds. This proves all of \eqref{eq:six}.

Finally index $(a,b,c,e)$ by $M=\{0,1,2,3\}$ and write
$H_i=G_{M\setminus\{i\}}$. The four groups $H_i$ are $R,Y,X,L$;
each has the intersection property. Equations~\eqref{eq:six} say precisely
that
\begin{equation}\label{eq:maximal}
 H_i\cap H_j=G_{M\setminus\{i,j\}}\qquad(i\ne j).
\end{equation}
Let $I,J\subseteq M$. If one is $M$, the desired intersection equality
is immediate. Otherwise choose $i\notin I$ and $j\notin J$.
If $i=j$, both subgroups lie in $H_i$, whose intersection property proves
the equality. If $i\ne j$, put $K=M\setminus\{i,j\}$.
An element $x\in G_I\cap G_J$ belongs to $H_i\cap H_j=G_K$.
Inside $H_i$, its membership in $G_I$ and $G_K$ gives
$x\in G_{I\cap K}$. Inside $H_j$, membership in this last subgroup and
$G_J$ gives
\[
 x\in G_{I\cap K\cap J}=G_{I\cap J},
\]
because $i\notin I$ and $j\notin J$. The opposite inclusion always
holds. Thus \eqref{eq:IP} holds for every pair $I,J$.
Together with the stipulated orders and \eqref{eq:string}, this is the
definition of the required string C-group.
\end{proof}

\begin{corollary}\label{cor:combined}
Suppose the distinguished involutions satisfy \eqref{eq:string}, the two
rank-three end subgroups satisfy the intersection property, and the point,
orbit, representatives and Schreier memberships of
Theorem~\ref{thm:ends} have been certified in the same group.
Then $(G;(a,b,c,e))$ is a string C-group of rank four.
\end{corollary}

\begin{proof}
Theorem~\ref{thm:ends} gives $L\cap R=D$; apply
Theorem~\ref{thm:rank4}.
\end{proof}

\section{What finite tables must establish}

For either rank-three end with distinguished index set $N$, it is enough
to certify $G_I\cap G_J=G_{I\cap J}$ for every pair of proper subsets
$I,J\subsetneq N$. Indeed, when one subset is $N$, the equality is
automatic. The proper subsets consist of the empty set, the three
singletons and the three pairs. Thus complete tables of the singleton and
pair-generated subgroups, with all their mutual intersections checked,
suffice; a table of the whole end subgroup is unnecessary.

Completeness is essential. A table $\mathcal T_I$ of elements is a complete
table of $G_I$ if it contains the identity, is closed under right
multiplication by every distinguished generator in $I$, and each listed
element is supplied as a word in those generators. Closure shows
$G_I\subseteq\mathcal T_I$, while the words show the reverse inclusion.
For involutions no additional inverse-generator closure is needed. The
empty table convention must instead be the one-element set
$\mathcal T_\varnothing=\{1\}$. Set comparisons between these complete
tables must use equality of elements in $G$.

The same completeness argument applies to the table of $D$. In the orbit
certificate, each $T_x$ must likewise be supplied as an actual word in
$a,b,c$, or its membership in $L$ must be established independently.
Knowing only that an ambient matrix carries $\ell$ to $x$ is insufficient.
Every point and every generator transition must be checked, with one
unambiguous representative for each point of the finite set. A truncated
orbit, a prefix of its edge list, or a proposed transversal without closure
does not satisfy Lemma~\ref{lem:schreier}.

\section{Projective points, matrix equality and orientation}

Let $V=\mathbb F_q^d$ consist of row vectors and let
$G\leq\GL(V)$ be the actual subgroup generated by the specified matrices.
Write $[v]$ for the line spanned by a nonzero row vector $v$, and define
$[v]^g=[vg]$. Matrix multiplication then agrees with the right-action
convention. If $x=[v]$ and $y=x^g$, the correct stabilizer element is
\[
 T_xgT_y^{-1}.
\]
Point equality may be tested projectively, for example by scaling the
first nonzero coordinate to one. This normalization changes the chosen
vector representing a point, not the group element $T_x$ representing the
path to that point. In particular a transversal matrix should be retained
as the actual product of its path generators. Multiplying it by an
arbitrary scalar to normalize its entries can destroy its proved
membership in $L$.

The action on projective points can have a nontrivial kernel. This causes
no defect in Lemma~\ref{lem:schreier}: kernel elements in $L$ are among
the stabilizer elements generated by the exact loops. However, when the
claimed group is the actual matrix subgroup $G\leq\GL(V)$, membership
of a loop in $D$ must be literal matrix membership. Equality to a table
entry only up to a nonzero scalar is not enough. Such a test proves at
most membership in the image of $D$ under the projective quotient. An
additional argument would be required to recover the assertion in $G$;
for example, it would suffice to establish independently that
$\ker(G\longrightarrow\PGL(V))\cap L\leq D$.

Alternatively, one may formulate all hypotheses and conclusions in the
projective image group from the outset. Then generator orders, word
identities, table intersections and subgroup membership must all refer
consistently to that image. This gives a statement about the image, not
automatically about an abstract group mapping onto it. A faithful
representation or a separately proved kernel argument is required for
such an additional transfer. No faithfulness of the projective point
action itself is required.

If an implementation instead uses column vectors with the direct left
action $x\mapsto gx$, and a representative satisfies $T_x\ell=x$,
then for $y=gx$ the stabilizer loop is
\[
 T_y^{-1}gT_x,
\]
not $T_xgT_y^{-1}$. These conventions must not be mixed. The right-action
formula can also be retained on column lines by defining
$x^g=g^{-1}x$ throughout. The fact that the generators themselves are
involutions does not make the two transversal conventions interchangeable:
a product representing $T_x$ need not be an involution.

The corollary proves the intersection property for the distinguished
matrix-generated group once its hypotheses have been verified. It does
not, by itself, identify that group with a named ambient group, prove its
order, enumerate all admissible markings, or complete a manuscript-level
classification. Those conclusions require their own independent evidence.

\section*{Attribution}

The telescoping argument in Lemma~\ref{lem:schreier} is the classical
Schreier generator argument; see Schreier~\cite{Schreier}. Theorem
\ref{thm:rank4} is the rank-four instance of the standard end-subgroup
criterion for string C-groups, conventionally cited as
McMullen--Schulte~\cite[Proposition 2E16]{MS}.
The complete proof above avoids any dependence on a classification of
finite simple groups or on a computational group-order routine.
The present contribution is only the explicit assembly of these
classical facts, with the action and certificate requirements stated for
the proposed application. No novelty is claimed for either underlying
criterion.

\begin{thebibliography}{9}
\bibitem{Schreier}
O.~Schreier, \emph{Die Untergruppen der freien Gruppen},
Abhandlungen aus dem Mathematischen Seminar der Universit\"at Hamburg
\textbf{5} (1927), 161--183.
\href{https://doi.org/10.1007/BF02952517}{doi:10.1007/BF02952517}.

\bibitem{MS}
P.~McMullen and E.~Schulte, \emph{Abstract Regular Polytopes},
Encyclopedia of Mathematics and its Applications, vol.~92,
Cambridge University Press, 2002, Chapter~2.
\href{https://doi.org/10.1017/CBO9780511546686.003}
{doi:10.1017/CBO9780511546686.003}.
\end{thebibliography}
\end{document}
